\documentclass{book}
\usepackage{amsmath}
\usepackage{amsthm}
\usepackage{amssymb}
\usepackage{tikz}
\usepackage{tikz-cd}
\usepackage{lipsum}
\usepackage{enumitem}

\DeclareMathOperator{\ouv}{Open}
\DeclareMathOperator{\Hom}{Hom}
\newtheorem{exam}{Example}
\title{Coherent Algeraic Sheaves}
\author{\textbf{Jean-Pierre Serre} \\ \\  Translated by Quentin Asparria and Axel Tamir}
\date{April 2026}
\newcounter{para}
\newcommand*{\numberedparagraph}{%
 \refstepcounter{para}\textbf{\thepara.}\space
}

\setlength{\parindent}{20pt}


\begin{document}

\maketitle

\tableofcontents

\chapter{Sheaves}

\section{Operations on Sheaves}

\subsection{Definition of a sheaf}

\par Let $X$ be a topological space. A \textit{sheaf of abelian groups} on $X$ (or simply a \textit{sheaf}) consists of:

\begin{enumerate}[label = (\alph*)]
	\item A \textit{function} $x \rightarrow \mathcal{F}_x$ \textit{sending every} $x \in X$ to an abelian group $\mathcal{F}_x$, 
	\item A topology on the set $\mathcal{F}$, the sum of the sets $\mathcal{F}_x$. 
\end{enumerate}

If $f$ is an element of $\mathcal{F}_x$, we set $\pi (f) = x$; the map $\pi$ is called the projection of $\mathcal F$ onto $X$; the subset of $\mathcal F \times \mathcal F$ formed by the pairs $(f,g)$ such that $\pi(f) = \pi(g)$ is denoted $\mathcal F + \mathcal F$. 

With these definitions in place, we can state the two axioms to which the data (a) and (b)  are subject:
\begin{enumerate}[label=(\Roman*)]

\item For all $f \in \mathcal F$, there exist open neighborhoods $V$ of $f$ and $U$ of $\pi(f)$ such that the restriction of $\pi$ at $V$ is a homeomorphism of $V$ onto $U$.
	(In other words, if $pi$ is a local homeomorphism)
\item The mapping $f \rightarrow -f$ is a continuous mapping of $\mathcal F$ into $\mathcal F$, and the mapping $(f,g) \rightarrow f + g$ is a continuous mapping of $\mathcal F + \mathcal F$ into $\mathcal F$. 

\end{enumerate}

It will be observed that, even if $X$ is separated (which we do not assume), $\mathcal F$ is not necessarily separated, which is illustrated by the example of germs of functions. 

\textbf{Example} of a sheaf. For an abelian group G, let us set $\mathcal F_x = G$ for every $x \in X$; the set $\mathcal F$ can be identified with the product $X \times G$, and, if we equip is with the product topology of $X$ by the discrete topology of $G$, we obtain a sheaf, called the \textit{constant sheaf} isomorphic with G, often identified with $G$.

\subsection{Sections of a sheaf}

Let $\mathcal F$ be a sheaf on the space $X$, and let $U$ be a subset of $X$. We define a \textit section of $\mathcal F$ over $U$ to be a continuous map $s: U \rightarrow \mathcal F$ such that $\pi \circ s$ is the identity mapping on $U$. So we have $s(x) \in \mathcal F_x$ for all $x \in U$. The set of sections of $F$ over $U$ is denoted by $\Gamma (U, \mathcal F)$; axiom (II) implies that $\Gamma (U, \mathcal F)$ is an abelian group. If $U \subset V$, and if $s$ is a section over $V$, the restriction of $s$ at $U$ is a section over $U$; hence we have a homomorphism $\rho_U^V: \Gamma (V, \mathcal F) \rightarrow \Gamma (U, \mathcal F)$. 

If $U$ is open on $X$, $s(U)$ is open in $\mathcal F$, and, $U$ runs over a base of the topology of $X$, then $s(U)$ is runs over a base of the topology of $\mathcal F$; this is only another way of stating axiom (I). 

Note another consequence of axiom (I): For every $f \in \mathcal F_x$, there exists a section $s$ over a neighborhood of $x$ such that $s(x) = f$, and two sections enjoying this property coincide on a neighborhood of $x$. In other words, $\mathcal F_x$ is an \textit{inductive limit} of $\Gamma (U, \mathcal F)$ following the filtering order of the neighborhoods of $x$. 

\subsection{Construction of sheaves}

Suppose we are given, for every open $U \subset X$, an abelian group $\mathcal F_U$, and, for every pair of open $U \subset V$, a homomorphism $\varphi_U^V: \mathcal F_V \rightarrow \mathcal F_V$, satisfying the transivity condition $\varphi_U^V \circ \varphi_V^W = \varphi_U^W$  whenever $U \subset V \subset W$. 

The collections $(\mathcal F_U, \varphi_U^V)$ allows us to define a sheaf in the following way:

\begin{enumerate}[label=(\alph*)]

	\item We set $\mathcal F_x = \lim \mathcal F_U$ (inductive limit with respect to the filtered ordering of open neighborhoods $U$ of $x$). If $x$ belongs to an open set $U$, we have a canonical morphism $\varphi_x^U: \mathcal F_U \rightarrow \mathcal F_x$. 

\item Let $t \in \mathcal F_U$, and denote by $[t,U]$ the set of $\varphi_x^U (t)$ for $x$ running over $U$; we have $[t,U] \subset \mathcal F$, and we give $\mathcal F$ the topology generated by $[t,U]$. Thus, an element $f \in \mathcal F_x$ admits as a neighborhood basis in $\mathcal F$ the sets $[t,U]$ where $x \in U$ and $\varphi_x^U (t) = f$. 

\end{enumerate}

We immediately verify that the data (a) and (b) satisfy the axioms (I) and (II), in other words, that $\mathcal F$ is indeed a sheaf. We say that this is the sheaf \textit{defined by the system} $(\mathcal F_U, \phi_U^V)$. 

If $t \in \mathcal F_U$, the mapping $x \rightarrow \phi_x^U (t)$ is a section of $\mathcal F$ over $U$; hence we have a canonical morphism $\iota: \mathcal F_U \rightarrow \Gamma (U, \mathcal F)$. 

\textbf{Proposition 1.} $\iota: \mathcal F_U \rightarrow \Gamma (U, \mathcal F)$ \textit{is injective if and only if the following condition holds}:

If an element $t \in \mathcal F_U$ has an open covering ${U_i}$ of $U$ with $\varphi_{U_i}^U (t) =0$ for all $i$, then $t=0$. 

If $t \in \mathcal F_U$ satisfies the previous condition, we have $$ \varphi_x^U (t) = \varphi_x^{U_i} \circ \varphi_{U_i}^U (t) =0 \text{if} x \in U_i, $$ which means that $\iota(t) =0$. Conversely, suppose that $\iota(t) =0$, with $t \in \mathcal F_U$; since $\varphi_x^U (t) =0$ for $x \in U$, there exists an open neighborhood $U(x)$ of $x$ such that $\varphi_{U(x)}^U (t) =0$, by the definition of an inductive limit. The sets $U(x)$ therefore form an open covering of $U$ satisfying the condition stated above. 

\end{document}
